\documentclass[../../../include/open-logic-section]{subfiles}

\begin{document}

\olfileid[fr]{sth}{choice}{countablechoice}
\olsection{Le choix dénombrable}

Sans recourir au choix ni au bon ordre,
il est facile de démontrer ceci :

\begin{lem}[dans $\Zminus$]
Tout ensemble fini admet une fonction de choix.
\end{lem}

\begin{proof}
Soit $a = \{b_1, \ldots, b_n\}$.
Supposons, pour simplifier, que chaque
$b_i \neq \emptyset$. Il existe alors
des objets $c_1, \ldots, c_n$ tels que
$c_1 \in b_1, \ldots, c_n \in b_n$.
D'après la
\olref[z][pairs]{prop:pairsconsequences},
l'ensemble
$\{\langle b_1, c_1\rangle , \ldots,
\langle b_n, c_n\rangle\}$ existe ;
c'est une fonction de choix pour~$a$.
\end{proof}

Mais dès que nous considérons des ensembles
infinis, les choses deviennent moins claires.
Prenons, par exemple, cette extension
«~minimale~» du résultat précédent :

\begin{defish}
\emph{Choix dénombrable.} Tout ensemble
\emph{dénombrable} admet une fonction de choix.
\end{defish}

Il s'agit d'un cas particulier de l'axiome
du choix. Or ce principe a été invoqué
assez souvent, sans que son emploi soit
toujours clairement perçu. En voici deux
exemples instructifs.\footnote{Ils sont
dus à \citet[\S9.4]{Potter2004}
et à Luca Incurvati.}

\begin{ex}
Voici une idée naturelle : pour tout ensemble
$A$, soit $\cardle{\omega}{A}$, soit
$\cardeq{A}{n}$ pour un certain
$n \in \omega$. C'est une façon
d'exprimer l'idée intuitive selon laquelle
tout ensemble est fini ou infini.
Cantor, comme bien d'autres mathématiciens,
a formulé cette affirmation sans la démontrer.
Pour notre part, nous l'avons démontrée
dans le
\olref[cardinals][classing]{generalinfinitycharacter}.
Mais nous y travaillions dans $\ZFC$ :
nous supposions que tout ensemble $A$
peut être bien ordonné et que
$\card{A}$ existe donc nécessairement.
Autrement dit, nous avions explicitement
supposé l'axiome du choix.

En fait, \citet{Dedekind1888} a proposé
sa propre démonstration de cette affirmation :

\begin{thm}[dans $\Zminus +
\text{choix dénombrable}$]
Pour tout $A$, soit
$\cardle{\omega}{A}$, soit
$\cardeq{A}{n}$ pour un certain
$n \in \omega$.
\end{thm}

\begin{proof}
Supposons que $\cardneq{A}{n}$ pour
tout $n \in \omega$. Alors, pour chaque
$n < \omega$, il existe en particulier
un sous-ensemble $A_n \subseteq A$
ayant exactement $\cardexpo{2}{n}$
éléments. À partir de la suite
$A_0, A_1, A_2, \ldots$, définissons
pour chaque $n$ :
\[
	B_n = A_n \setminus
	\bigcup_{i < n} A_i.
\]
Remarquons que :
\begin{align*}
	\card{\bigcup_{i < n}A_i}
	&\leq \card{A_0} + \card{A_1}
	+ \ldots + \card{A_{n-1}}\\
	&=1 + 2 + \ldots + 2^{n-1}\\
	& = 2^n - 1\\
	& < 2^n = \card{A_n}
\end{align*}
\footnote{Note de l'édition française :
l'original écrit $A_n$ sous le signe
d'union dans la première ligne ;
la borne porte sur l'union des $A_i$
pour $i<n$.}
Chaque $B_n$ possède donc au moins
un élément $c_n$. De plus, les $B_n$
sont deux à deux disjoints : si
$c_n=c_m$, alors $n=m$.
Or tous les $c_n$ appartiennent à $A$.
La fonction $f(n)=c_n$ est donc une
injection de $\omega$ dans $A$.
\end{proof}
\noindent
Dedekind n'a pas signalé qu'il avait
utilisé le choix dénombrable.
Mais \emph{vous}, avez-vous repéré
où il intervient ? Relisez la
démonstration. (Vraiment :
\emph{relisez-la}.)

Elle utilise deux fois le choix
dénombrable : d'abord pour obtenir
la suite d'ensembles
$A_0, A_1, A_2, \ldots$,
puis pour choisir un élément $c_n$
dans chaque~$B_n$. Cet usage du choix
est en outre indispensable.
\citet[p.~138]{Cohen1966} a montré
que le résultat peut échouer en l'absence
de toute forme de choix : il est
compatible avec $\ZF$ qu'il existe
des ensembles \emph{incomparables}
à~$\omega$.
\end{ex}

\begin{ex}
En \citeyear{Cantor1878}, Cantor a affirmé
qu'une union dénombrable d'ensembles
dénombrables est dénombrable.
Il n'a pas donné de démonstration,
ce qui suggère peut-être qu'elle
lui paraissait évidente. Pour notre
part, nous avons démontré une version
plus générale de ce résultat dans la
\olref[card-arithmetic][simp]{kappaunionkappasize}.
Mais notre démonstration supposait
explicitement l'axiome du choix.
Même la démonstration du résultat moins général
fait appel au choix dénombrable.

\begin{thm}[dans $\Zminus +
\text{choix dénombrable}$]
Si $A_n$ est dénombrable pour chaque
$n \in \omega$, alors
$\bigcup_{n < \omega} A_n$
est dénombrable.
\end{thm}

\begin{proof}
Sans perte de généralité, supposons
que chaque $A_n \neq \emptyset$.
Pour chaque $n \in \omega$, il existe
alors !!a{surjection}
$f_n \colon \omega \to A_n$.
Définissons
$f \colon \omega \times \omega
\to \bigcup_{n < \omega} A_n$
par $f(m,n)=f_n(m)$.
Le résultat suit du fait que
$\omega \times \omega$ est
dénombrable
(\olref[sfr][siz][zigzag]{natsquaredenumerable})
et que $f$ est !!a{surjection}.
\end{proof}
\noindent
Avez-vous repéré l'emploi du choix
dénombrable ? Il sert à choisir
la suite de fonctions
$f_0, f_1, f_2, \dots$\footnote{
Un usage semblable du choix figurait
dans la
\olref[card-arithmetic][simp]{kappaunionkappasize},
lorsque nous avons demandé :
«~Pour chaque
$\beta \in \cardfont{a}$,
fixons !!a{injection} $f_\beta$~».}
Là encore, le résultat peut échouer
en l'absence de tout principe de choix.
Plus précisément,
\citet{FefermanLevy1963} ont démontré
qu'il est compatible avec $\ZF$
qu'une union dénombrable d'ensembles
dénombrables soit équipotente à
$\Pow{\omega}$.\footnote{Note de l'édition
française : l'original exprime cette taille
par $\beth_1$. Sans l'axiome du choix,
$\Pow{\omega}$ peut ne pas être
bien ordonnable, et la définition
antérieure de $\beth_1$ comme cardinal
ordinal ne s'applique alors pas.
L'énoncé d'équipotence évite cette
présupposition.}
Mais Russell a exprimé cette idée
de manière bien plus amusante :
\begin{quote}
	C'est ce qu'illustre le millionnaire
	qui achetait une paire de chaussettes
	chaque fois qu'il achetait une paire
	de bottes, et jamais à une autre
	occasion. Il avait une telle passion
	pour ces achats qu'il finit par
	posséder $\aleph_0$ paires de bottes
	et $\aleph_0$ paires de chaussettes
	\dots\@ Parmi les bottes, nous
	pouvons distinguer la droite de la
	gauche et choisir ainsi une botte
	dans chaque paire : toutes les bottes
	droites, ou toutes les gauches.
	Avec les chaussettes, aucun principe
	de sélection semblable ne se présente
	à nous. Sans supposer l'axiome
	multiplicatif [c'est-à-dire,
	en substance, l'axiome du choix],
	nous ne pouvons être certains
	qu'il existe une classe contenant
	une chaussette de chaque paire.
	\citep[p.~126]{Russell1919}
\end{quote}
Autrement dit, une certaine forme
de choix est nécessaire pour
démontrer que, si l'on possède
une infinité dénombrable de paires
de chaussettes, on possède un ensemble
dénombrable de chaussettes.
En effet, sans le choix dénombrable
(ou un principe équivalent), une union
dénombrable d'ensembles dénombrables
peut ne pas être dénombrable.
\end{ex}

La leçon est que le choix dénombrable
a été utilisé à maintes reprises
sans que ses utilisateurs en aient
toujours conscience. La question
philosophique est de savoir comment
nous pourrions le \emph{justifier}.

Une tentative de justification
intuitive pourrait invoquer une
supertâche. Supposons que nous
fassions le premier choix en
$\nicefrac{1}{2}$ minute,
le deuxième en $\nicefrac{1}{4}$
minute, \dots, le $n$-ième en
$\nicefrac{1}{2^n}$ minute,
\dots\@ En $1$~minute, nous aurions
effectué une suite de choix de
longueur $\omega$ et défini une
fonction de choix.

Mais que pourrait réellement nous
apprendre cette expérience de pensée ?
D'abord, elle prend l'idée de
«~choisir~» au pied de la lettre.
Ensuite, elle semble lier les
mathématiques à la possibilité
métaphysique.

Surtout, elle ne justifie nullement
le choix \emph{tout court}, au-delà
du \emph{seul} choix dénombrable.
Car, si \emph{tout} ensemble doit
admettre une fonction de choix,
il nous faudrait pouvoir accomplir
une «~supertâche de longueur ordinale
arbitraire~». Franchement, cette
idée est risible.

\end{document}
